\section{Sensor Calibration problem}

\subsection{Problem Description}
The objective of this problem is to determine the coefficients ($A, B, C$) of the Steinhart-Hart equation to model the relationship between the resistance ($R$) and temperature ($T$) of a thermistor. The constitutive equation is defined as:
\begin{equation}
\frac{1}{T} = A + B \ln(R) + C (\ln(R))^3
\end{equation}
Where $T$ is in Kelvin and $R$ is in Ohms. The optimisation seeks to minimise the Mean Squared Error (MSE) between the model-predicted temperature and the measured data points. Unlike the Truss problem, this problem is unconstrained but highly non-linear and ill-conditioned due to the varying orders of magnitude of the coefficients.

\subsection{Method}
The BFGS algorithm, implemented through `scipy.optimize.minimize` with the method 'L-BFGS-B' was used to minimise the MSE. A key aspect of this approach was utilising the analytical gradient (Jacobian) to steer the solver, rather than relying on finite-difference approximations.
To thoroughly compare the robustness of the deterministic BFGS method with the stochastic Simulated Annealing (SA) method from Task 3, three specific protocols were carried out.
To thoroughly compare the robustness of the deterministic BFGS method with the stochastic Simulated Annealing (SA) method from Task 3, three specific protocols were designed to simulate different types of operational uncertainties.

\textbf{Protocol A: Measurement Noise (Realistic Sensor Error)}
This protocol simulates the most common real-world scenario: inaccuracies in the measuring equipment. Gaussian white noise is added to the temperature measurements ($T_{meas}$) \textit{before} the optimization begins.
\begin{equation}
T_{meas, i} = T_{true, i} + \mathcal{N}(0, \sigma^2)
\end{equation}
Critically, this noise is static during the optimization process. This means the objective function landscape remains smooth and differentiable, as the "target" data points are fixed, even if they are slightly wrong. This tests the solver's ability to find the best fit through noisy data (Least Squares regression).

\textbf{Protocol B: Function Noise (Stochastic Evaluation)}
This protocol simulates a more hostile environment where the objective function evaluation itself is unstable. A random noise term is added to the Mean Squared Error (MSE) calculation at \textit{every} step of the optimization.
\begin{equation}
f(x) = MSE(x) + \mathcal{N}(0, \sigma^2)
\end{equation}
This represents scenarios such as Monte Carlo-based evaluations or systems with internal fluctuations (e.g., analog voltage jitter during readout). Unlike Protocol A, this makes the objective function non-deterministic: evaluating the same point twice yields different results, creating a "shuddering" landscape that challenges the gradient estimation.

\textbf{Protocol C: Robustness Stress Test}
This protocol determines the breaking point of the algorithm. By gradually increasing the noise intensity ($\sigma$) from negligible levels ($10^{-4}$) to extreme levels ($5.0$), we construct a "Robustness Curve". This allows us to identify the specific noise threshold where the precision of the deterministic BFGS method degrades below that of the stochastic Simulated Annealing method, marking the "crossover point" where switching algorithms becomes necessary.

\subsection{Results}

\subsubsection{Exploration of Performance (Protocol A \& B)}
Protocol A tested the solver against measurement noise (Gaussian noise added to $T_{data}$).

\textbf{Quantitative Analysis of Protocol A Results}

The table below summarises the data extracted from the log file. It compares Single-Start BFGS (one run per seed) to Multi-Start BFGS (best of 5 runs per seed) across 500 different noise seeds.

\begin{figure}[H]
    \centering
    \begin{subfigure}[b]{0.48\textwidth}
        \includegraphics[width=\textwidth]{protocol_a_seed0_single_objective.png}
        \caption{Objective Function}
    \end{subfigure}
    \begin{subfigure}[b]{0.48\textwidth}
        \includegraphics[width=\textwidth]{protocol_a_seed0_single_gradient.png}
        \caption{Gradient Norm}
    \end{subfigure}
    \begin{subfigure}[b]{0.48\textwidth}
        \includegraphics[width=\textwidth]{protocol_a_seed0_single_step_size.png}
        \caption{Step Size}
    \end{subfigure}
    \begin{subfigure}[b]{0.48\textwidth}
        \includegraphics[width=\textwidth]{protocol_a_seed0_single_model_fit.png}
        \caption{Model Fit}
    \end{subfigure}
    \caption{Protocol A: Detailed Convergence Analysis}
    \label{fig:sensor_proto_a_detailed}
\end{figure}

\begin{table}[H]
    \centering
    \resizebox{\textwidth}{!}{
    \begin{tabular}{|l|l|l|l|}
    \hline
    Metric & Single-Start BFGS & Multi-Start BFGS (Best of 5) & Improvement \\ \hline
    Success Rate & 100\% (500/500 converged) & 100\% (500/500 converged) & +0\% \\ \hline
    Mean Final MSE & 0.2314 & 0.2314 & $\approx$ 0.0\% (Identical) \\ \hline
    Avg. Function Evals & 38.4 & 183.2 (Total for 5 starts) & N/A \\ \hline
    Avg. Execution Time & 0.009s & 0.045s & 5x Slower \\ \hline
    \end{tabular}
    }
    \caption{Protocol A Comparison}
    \label{tab:sensor_protocol_a}
\end{table}

Multi-Start BFGS took 5 times longer (0.045s vs. 0.009s) but did not improve accuracy.
\textbf{Engineering conclusion}: For this particular type of problem (sensor calibration with static measurement noise), using Multi-Start is inefficient. A single deterministic run is both sufficient and optimal.

Adding Gaussian noise to the measurements ($T_{data}$) does not compromise the smoothness of the objective function. The average Final MSE for Single-Start (0.2314) matches that of Multi-Start (0.2314) up to four decimal places. If the objective function landscape were non-convex with many local minima, the Single-Start method would often get trapped in a poor valley, resulting in a higher average MSE. Its identical performance with Multi-Start suggests that, for the Steinhart-Hart equation, the optimisation landscape is effectively unimodal (similar to convex) even when noisy data is involved.

\textbf{Efficiency of Deterministic Gradients}: BFGS solved the problem in an average of 38 function evaluations, which is significantly faster compared to the Task 3 Simulated Annealing results, typically requiring over 3,000 evaluations. Since the noise is incorporated into the data before the solver begins, the gradient $\nabla f(x)$ remains consistent. This enables the BFGS algorithm to construct a precise approximation of the Hessian (curvature), allowing it to take "Newton-like" quadratic steps rather than only small downhill moves.

Protocol B tested the solver against function noise (noise added to the MSE calculation). The convergence plots reveal a rapid failure mode.

\begin{table}[H]
    \centering
    \begin{tabular}{|l|l|l|l|}
    \hline
    Metric & Protocol A & Protocol B & Status \\ \hline
    Stability & Rock solid (MSE 0.23) & Chaotic (MSE -0.0001 to 71.0) & FAILED \\ \hline
    Reliability & 100\% Consistent & High Variance & FAILED \\ \hline
    Avg. Evals & $\approx$ 38 & $\approx$ 25 (Premature Stop) & FAILED \\ \hline
    Behavior & Found global minimum & Random noise spikes & FAILED \\ \hline
    \end{tabular}
    \caption{Protocol A vs B Comparison}
    \label{tab:sensor_protocol_ab}
\end{table}

BFGS encounters premature convergence issues in Protocol B. It often becomes trapped in 'phantom minima' caused by noise spikes, concluding after as few as 12 iterations with unacceptable error levels.

\begin{figure}[H]
    \centering
    \begin{subfigure}[b]{0.48\textwidth}
        \includegraphics[width=\textwidth]{protocol_b_seed0_single_objective.png}
        \caption{Objective Function}
    \end{subfigure}
    \begin{subfigure}[b]{0.48\textwidth}
        \includegraphics[width=\textwidth]{protocol_b_seed0_single_gradient.png}
        \caption{Gradient Norm}
    \end{subfigure}
    \begin{subfigure}[b]{0.48\textwidth}
        \includegraphics[width=\textwidth]{protocol_b_seed0_single_step_size.png}
        \caption{Step Size}
    \end{subfigure}
    \begin{subfigure}[b]{0.48\textwidth}
        \includegraphics[width=\textwidth]{protocol_b_seed0_single_model_fit.png}
        \caption{Model Fit}
    \end{subfigure}
    \caption{Protocol B: Detailed Convergence Analysis}
    \label{fig:sensor_proto_b_detailed}
\end{figure}

The problem of Multi-Start in Protocol B is evident. In Protocol A, Multi-Start was wasteful but consistent. In Protocol B, Multi-Start becomes random. Simulated Annealing (Task 3) would handle this better because the Metropolis criterion ($P = e^{-\Delta E / T}$) allows it to ignore small upward noise spikes and continue searching, effectively "smoothing" the landscape. BFGS has no such mechanism.

\subsubsection{Exploration of Robustness (Protocol C)}
The Robustness Curve (Figure \ref{fig:sensor_robustness}) compares the final MSE of Single BFGS, Multi-Start BFGS, Simulated Annealing (SA), and a Hybrid (SA+BFGS) approach across increasing noise levels ($\sigma$).

\begin{figure}[H]
    \centering
    \includegraphics[width=0.8\textwidth]{robustness_curve.png}
    \caption{Robustness Curve}
    \label{fig:sensor_robustness}
\end{figure}

Observation:
\begin{itemize}
    \item BFGS: The error increases linearly (on a log-log scale) with the noise level $\sigma$.
    \item Simulated Annealing: The SA error remains effectively constant and high (forming a flat line), regardless of the noise level.
    \item No cross over clearly visible with pure SA.
\end{itemize}

\begin{figure}[H]
    \centering
    \begin{subfigure}[b]{0.48\textwidth}
        \includegraphics[width=\textwidth]{protocol_c_sigma0.01_run0_objective.png}
        \caption{Objective Function}
    \end{subfigure}
    \begin{subfigure}[b]{0.48\textwidth}
        \includegraphics[width=\textwidth]{protocol_c_sigma0.01_run0_gradient.png}
        \caption{Gradient Norm}
    \end{subfigure}
    \begin{subfigure}[b]{0.48\textwidth}
        \includegraphics[width=\textwidth]{protocol_c_sigma0.01_run0_step_size.png}
        \caption{Step Size}
    \end{subfigure}
    \begin{subfigure}[b]{0.48\textwidth}
        \includegraphics[width=\textwidth]{protocol_c_sigma0.01_run0_model_fit.png}
        \caption{Model Fit}
    \end{subfigure}
    \caption{Protocol C: Low Noise Regime ($\sigma=0.01$) Analysis}
    \label{fig:sensor_proto_c_detailed}
\end{figure}

\begin{table}[H]
    \centering
    \begin{tabular}{|l|l|l|l|l|}
    \hline
    Noise Level ($\sigma$) & BFGS (Single) & SA (Task 3) & Hybrid (SA+BFGS) & Winner \\ \hline
    0.01 (Low) & 9.50e-5 & 1.39e-4 & 9.50e-5 & BFGS / Hybrid \\ \hline
    0.05 (Medium) & 0.00214 & 0.00224 & 0.00214 & BFGS / Hybrid \\ \hline
    0.10 (High) & 0.00855 & 0.00858 & 0.00855 & Tie \\ \hline
    \end{tabular}
    \caption{Noise Level Performance Comparison}
    \label{tab:sensor_noise_perf}
\end{table}

At low noise ($\sigma=0.01$), Simulated Annealing is roughly 46\% worse than BFGS. This confirms that SA struggles when errors start reducing. Even at low temperatures, the probabilistic nature of SA prevents it from locking onto the exact mathematical minimum.
The Hybrid method matches the precision of BFGS exactly at low noise. This proves the Hybrid implementation works as intended:
Phase 1 (SA): Finds the correct valley (Global Search).
Phase 2 (BFGS): Takes the SA result and performs a gradient descent "polish" to hit the exact bottom (Local Search).

\subsubsection{The Performance Profile (Dolan-More)}

\begin{figure}[H]
    \centering
    \includegraphics[width=0.8\textwidth]{performance_profile_A.png}
    \caption{Performance Profile}
    \label{fig:sensor_perf_prof}
\end{figure}

It compares Probability of Success against Computational Cost. The Robustness Curve reveals a distinct 'crossover point' at $\sigma \approx 0.1$. Below this threshold, the BFGS component of the Hybrid solver drives the error to machine precision however, its reliability degrades in the presence of function noise (Protocol B).

\subsection{Discussion: Algorithmic Sensitivity \& Robustness}

\subsubsection{Computational Efficiency vs. Redundancy (Protocol A)}
The results from Protocol A provide a critical insight into the utility of Multi-Start strategies in gradient-based optimization.

\textbf{The Cost of Safety}: As quantified by the Dolan-Moré performance profile \cite{dolan2002}, the Multi-Start approach exhibited a performance ratio of $\tau \approx 5$. This accurately reflects the computational burden; the algorithm evaluated five independent trajectories to locate the same minimum that the Single-Start method found in one.

\textbf{The "Unimodal" Conclusion}: Since the Single-Start method achieved a 100\% convergence rate with an MSE identical to the Multi-Start method (0.2314), we can conclude that measurement noise does not create significant local minima in the Steinhart-Hart landscape.

\textbf{Engineering Decision}: This finding contradicts the general recommendation in global optimization literature \cite{rao2019} to always employ restarts for non-linear problems. For this specific class of regression problem, Multi-Start was experimentally proven to be computationally redundant.

\subsubsection{Robustness to Noise Topologies}
The divergence in performance between Protocols A and B highlights the fundamental dependence of Quasi-Newton methods on the validity of the gradient approximation.

\textbf{A. Measurement Noise (Protocol A): The Smoothing Effect}
BFGS demonstrated high robustness to measurement noise. This aligns with the theory of Least Squares estimation.
\textbf{Mechanism}: The objective function $f(x)$ calculates the Mean Squared Error over $N$ data points. This summation operation effectively acts as a low-pass filter, averaging out the Gaussian noise ($\epsilon \sim N(0, \sigma^2)$) attached to the measurements \cite{bishop2006}.
\textbf{Result}: Consequently, the underlying topological landscape remains continuous and differentiable, allowing the L-BFGS-B algorithm to construct a valid Hessian approximation and achieve quadratic convergence rates (avg. 38 evaluations).

\textbf{B. Function Noise (Protocol B): The Wolfe Condition Failure}
In contrast, BFGS failed catastrophically when noise was injected directly into the function value (Protocol B). This confirms the theoretical limitations of deterministic solvers in stochastic environments as described by Nocedal \& Wright \cite{nocedalwright2006}.
\textbf{Mechanism of Failure}: The Random Jitter in $f(x)$ prevents the Line Search algorithm from satisfying the Wolfe Conditions (specifically the Armijo Rule or "Sufficient Decrease" condition).
\textbf{Consequence}: As seen in Seed 96, the solver interpreted transient noise spikes as physical curvature. This corrupted the BFGS update formula, causing the Hessian approximation to become ill-conditioned. The solver was forced to shrink the step size to machine epsilon, triggering the "premature termination" observed (avg. 12 evaluations).

\textbf{Comparison to Task 3}: This failure stands in sharp contrast to the Simulated Annealing results from our Task 3 Report \cite{task3report}. SA’s Metropolis criterion ($P = e^{-\Delta E / T}$) allows for uphill moves, effectively "smoothing" these noise spikes, whereas BFGS treats every fluctuation as a rigid mathematical constraint.
